\documentclass[aspectratio=169]{beamer}
\usepackage{tikz}
\usetikzlibrary{shadings}
\usepackage[T2A]{fontenc} % ДОБАВЬТЕ ЭТУ СТРОКУ
\usepackage[utf8]{inputenc}
\usepackage{amsmath,amssymb}
\usepackage[russian]{babel}
\usepackage{xcolor}
\usepackage{graphicx}
\author[Ежелев Г. И.]{Ежелев Г. И.}

\newlength{\RightSidebarWidth}
\setlength{\RightSidebarWidth}{0.8cm}

% используем внешний шаблон sidebar
\useoutertheme[right]{sidebar}

% убираем стандартный верхний и нижний колонтитулы
\setbeamertemplate{headline}{}
\setbeamertemplate{footline}{}

% наш вариант правой боковой панели
\setbeamertemplate{sidebar right}{%
	\begin{tikzpicture}[remember picture, overlay]
		% Можно использовать structure.fg напрямую
		\shade[top color=structure.fg!60, bottom color=structure.fg!10]
		([xshift=-2.15cm]current page.north east) rectangle (current page.south east);
		\node[anchor=center, align=center, font=\tiny, text width=2cm]
		at ([xshift=-1cm, yshift=-1cm]current page.north east) {  % xshift чуть левее
			\footnotesize \textcolor{structure.fg}{\textbf{STM32}}\\[0.3em]
			\tiny \myshorttitle\\[0.3em]
			\tiny \textit{\insertauthor}
		};
	\end{tikzpicture}
}

% чтобы содержимое не залезало под полосу
\setbeamersize{text margin right=\RightSidebarWidth}
\usepackage{newunicodechar}
\usepackage{wasysym}
\newunicodechar{😋}{\smiley}
\def\myshorttitle{Глава 12. \\ Многопоточность. RTOS}
\title[\myshorttitle]{Глава 12. Многопоточность. RTOS}

\begin{document}
	\begin{frame}
	\includegraphics[width=10cm, height=8cm]{"mem.jpeg"}
\end{frame}
	\titlepage
	\newpage
	\large \textbf{\textcolor{structure.fg}{Управление задачами}}\\
	$\bullet$ Весь код, который вы писали до этого - выполнялся последовательно. То есть был какой-то \textcolor{orange}{основной цикл} в котором выполнялся весь код.\\
	$\bullet$ Такой подход называется \textcolor{structure.fg}{Super loop}\\
	\begin{itemize}
		\item Такой код быстрее писать и удобнее понимать
	\item Завис один узел - \textcolor{orange}{гарантированно зависла вся программа}
	\item Сложно писать распределение нагрузки. Если регулятор нужно обновлять раз в 100 мс, раз в 10 мс нужно опрашивать датчик и раз в 30 мс отправлять пакеты экрану - на каждое действие должен быть свой таймер (а если частота важна - то прерывание).
	\item Многозадачность реализуется \textcolor{structure.fg}{конечным автоматом}.
\end{itemize}
\includegraphics[width=6cm, height=3cm]{rtos2.jpg}
\newpage
\large \textbf{\textcolor{structure.fg}{Конечный автомат}}\\
\textcolor{structure.fg}{Конечный автомат} (\textcolor{orange}{Finite State Machine, FSM}) - это программа, имеющая конечное число состояний и правил перехода между ними. В целом довольно простая мысль. Суть в том, как эти состояния реализовать.\\
В случае робототехники:\\
\begin{itemize}
	\item \textcolor{orange}{Цикл делится на 3 части}: \textcolor{structure.fg}{Listen-Think-Act} (Слушай-Думай-Делай).
	\item \textcolor{structure.fg}{Listen} - первая часть (в самом начале), занимается обработкой датчиков (\textcolor{red}{никаких чтений датчиков в других частях цикла - это костыль})
	\item \textcolor{structure.fg}{Think} - Конечный автомат. Он обычно делится на две части.\\$\bullet$ \textcolor{orange}{Первая} - только выбор состояния на основании данных с датчиков. То есть набор условий:\\
	if(state == 1) {[условия перехода из 1 -> \{2,..., n\}]} \\.... \\if(state == n){[условия перехода из 1 -> \{1,..., n-1\}]}
	\newpage
	$\bullet$ \textcolor{orange}{Вторая часть} - только выполнение состояний (без переходов между ними):\\
	if(state == 1) {[Действие в состоянии 1]} \\.... \\if(state == n){[Действие в состоянии n]}
	$\bullet$ Такой подход позволяет \textcolor{orange}{не запутаться во всех состояниях}, когда код разрастеться.
	\item \textcolor{structure.fg}{Act} - В конце - отдача приказов узлам робота (моторы, экранчик и т.п.)\textcolor{red}{никаких моторов в середине кода!}.
	\item \textcolor{structure.fg}{Listen-Think-Act} - это пример паттерна проектирования - это некоторый шаблон или идея, как надо реализовывать какой-то кусок кода (т.к. не вы первые догадались датчики читать и потом по данным с них двигаться).
	\item Многие паттерны звучат очевидно и скучно, но при этом почему-то люди всё равно их игнорируют и пишут плохой код {\huge\smiley}
\end{itemize}
\newpage
\large \textbf{\textcolor{structure.fg}{Пример. Парсинг строки}}\\
\includegraphics[width=12cm, height=8cm]{rtos1.png}
\newpage
\large \textbf{\textcolor{structure.fg}{Параллельные задачи}}\\
\includegraphics[width=13cm, height=8cm]{rtos3.png}
\newpage
\large \textbf{\textcolor{structure.fg}{Параллельные задачи. Блокировки}}\\
\includegraphics[width=13cm, height=8cm]{rtos4.png}
\large \textbf{\textcolor{structure.fg}{Параллельные задачи. Переключение контекста}}\\
\includegraphics[width=13cm, height=8cm]{rtos5.png}
\large \textbf{\textcolor{structure.fg}{Параллельные задачи. Конкурентность}}\\
\includegraphics[width=13cm, height=8cm]{rtos6.png}
\newpage
\large \textbf{\textcolor{structure.fg}{Конкурентность}}\\
\begin{itemize}
	\item \textcolor{structure.fg}{Конкурентность} - исполнение нескольких задач \textcolor{orange}{одновременно}. Это не всегда значит, что задачи физических работают одновременно. \textcolor{orange}{Главное - они реализованы, как отдельные независимые задачи}, которые не знают о существовании других.
	\item В случае, если одна задача отдает на время управление, \textcolor{orange}{программа переключается на исполнение другой задачи}. В этот момент происходит \textcolor{structure.fg}{переключение контекста} - сохранение данных одной задачи и подгрузка данных следующей.
	\item \textcolor{orange}{Даже с одним ядром - такой подход имеет свои плюсы} - вам не нужно думать, как распределять ресурсы между задачами (моторы раз в 100 мс, регуляторы раз в 10...). Каждая идейная задача - отдельная параллельная задача в коде, система сама переключит контекст когда надо.
\end{itemize}
\newpage
\large \textbf{\textcolor{structure.fg}{Итог}}\\
\includegraphics[width=12cm, height=8cm]{rtos7.png}
\newpage
\large \textbf{\textcolor{structure.fg}{Варианты архитектур}}\\
\includegraphics[width=12cm, height=8cm]{rtos8.png}
\newpage 
\large \textbf{\textcolor{structure.fg}{RTOS}}\\
Все \textcolor{orange}{манипуляции с процессами} (например переключение контекста) осуществляет \textcolor{structure.fg}{планировщик(Scheduler)} - часть операционной системы.\\
$\bullet$\textcolor{structure.fg}{Операционные системы реального времени} (Real Time Operating System, RTOS) - это ОС, которые ставят в приоритет не равномерность распределения ресурсов между задачами и удобство пользователя, а \textcolor{orange}{временную детерминированность}, т.е. реакция на событие произойдет в строго(почти) определенное время.\\
Пример такой системы - FREERTOS, состоящей из:
\begin{itemize}
	\item \textcolor{structure.fg}{Планировщика(scheduler)} - решает, какая задача выполняется сейчас
	\item \textcolor{structure.fg}{Диспетчера задач (Task manager)} - следит за состоянием задачи, памятью и тп.
	\item \textcolor{structure.fg}{Управление памятью (Heap manager)} - управление данными задач (их целых 4-5 штук)
	\item \textcolor{structure.fg}{Механизмы синхронизации} - очереди, семафоры, уведомления, группы событий.
\end{itemize}
\newpage
\large \textbf{\textcolor{structure.fg}{Популярность RTOS}}\\
По опросу Eclipse Foundation для IoT (Internet of Things) устройств после Linux и Windows- \textcolor{orange}{самой популярной основой стал FREERTOS}, обогнав разработку без операционной системы (Bare Metal)\\
\includegraphics[width=10cm, height=6.5cm]{linux.png}
\newpage
\large \textbf{\textcolor{structure.fg}{Задачи во FREERTOS}}\\
\textcolor{structure.fg}{Задача} - это какая-то функция, обычно реализованная как (почти)бесконечный цикл, планировщик их сам нарезает:\\
\begin{tikzpicture}[remember picture,overlay]
	
	\node[anchor=north west] at ([xshift=-1cm,yshift=-1.5cm]current page.north west) {
		\begin{tabular}{c}
			\includegraphics[width=8cm, height=6cm]{task1.png}\\
			{Как мы видим задачи в коде}
		\end{tabular}
	};
	
	\node[anchor=north west] at ([xshift=7cm,yshift=-1.5cm]current page.north west) {
		\begin{tabular}{c}
			\includegraphics[width=8cm, height=6cm]{task2.png}\\
			{Их реальное исполнение}
		\end{tabular}
	};
\end{tikzpicture}
\newpage
\large \textbf{\textcolor{structure.fg}{Переключение задач}}\\
\begin{itemize}
	\item \textcolor{orange}{Каждую миллисекунду FreeRtos пробуждает ядро по тому самому прерыванию SysTick}. Каждый период между пробуждениями ОС так и называется - "тик".
	\item Задачи \textcolor{structure.fg}{переключаются операционкой в следующих ситуациях}:
	\item \textcolor{orange}{По приоритету задач}. В новый тик есть заждавшаяся задача того же приоритета или любая задача более высокого приоритета. Это реализовано через массив очередей(у каждого уровня приоритетов своя очередь) - pxReadyTasksLists[configMAX\_PRIORITIES].
	\item \textcolor{orange}{Вытеснение (Preemption)}. Если неважная задача А ставиться на паузу прерыванием, например приемом пакета UART (а прерывания никто не отменял), а после приема пакета в очередь попадает более важная задача Б, то RTOS сразу последует в Б, не завершая А.
	\newpage
	\item \textcolor{orange}{Ожидание}. Мы можем поставить задачу в ожидание, пока какая-то другая задача что-то не сделает, что пробудит нашу задачу.
	\item \textcolor{orange}{Блокировка}. Задачи могут обмениваться данными. Если задача попробует запросить данные у другой, которых еще нет, она заблокируется, пока данные не появятся, а в это время будут исполняться другие таски
\end{itemize}
\newpage
\large \textbf{\textcolor{structure.fg}{Состояния задач}}\\
\includegraphics[width=9cm, height=7.5cm]{task3.png}\\
\newpage
\large \textbf{\textcolor{structure.fg}{Состояния задач}}\\
\begin{itemize}
	\item \textcolor{structure.fg}{Ready} - После создания задача попадает в это состояние. Такие задачи готовы к исполнению и ждут внимания ядра
	\item \textcolor{structure.fg}{Running} - выполнение
	\item \textcolor{structure.fg}{Blocked} - заблокирована. Например вызвана задержка, или задача попыталась запросить данные от другой задачи, а их нет (Блокировка) или задача ожидает событие (Ожидание). В этом состоянии задача лежит в планировщике.
	\item \textcolor{structure.fg}{Suspended} - приостановлена. Более радикальная мера, задача выбрасывается из планировщика. Нужна для отключения какого-то куска логики. Зачем нужно? => чтобы сохранить контекст и данные самой задачи.
\end{itemize}
\newpage
\large \textbf{\textcolor{structure.fg}{сборка проекта FREERTOS}}\\
Все файлы операционки уже есть в проекте. проверьте что все нужные файлы добавлены в keil:\\[5cm]
\begin{tikzpicture}[remember picture,overlay]
	
	\node[anchor=north west] at ([xshift=0cm,yshift=-1.5cm]current page.north west) {
		\begin{tabular}{c}
			\includegraphics[width=5cm, height=4cm]{keil1.png}\\
			{FreeRtosConfig.h}
		\end{tabular}
	};
	
	\node[anchor=north west] at ([xshift=10.5cm,yshift=-1.5cm]current page.north west) {
		\begin{tabular}{c}
			\includegraphics[width=5cm, height=4cm]{keil2.png}\\
			{Все файлы операционки}
		\end{tabular}
	};
		\node[anchor=north west] at ([xshift=5cm,yshift=-1.5cm]current page.north west) {
		\begin{tabular}{c}
			\includegraphics[width=5cm, height=4cm]{keil3.png}\\
			{Не забудьте подключить}
		\end{tabular}
	};
\end{tikzpicture}
\textcolor{red}{Закоментировать SysTick в нашем Time.h, FREERTOS его тоже использует для синхронизации}\\
Советую посмотреть файл FreeRtosConfig.h (он выложен в репе, скопируйте к себе, в проекте он возможно неправильный)
\newpage
\large \textbf{\textcolor{structure.fg}{Создание задачи}}\\
\includegraphics[width=6cm, height=5cm]{task4.png}\\
	\begin{tikzpicture}[remember picture,overlay]
	\node[draw,rounded corners,fill=green!3,text width=9cm,align=left]
	at (10.5,2.25) 
	{\begin{itemize}
			\item \textcolor{structure.fg}{pvTaskCode} - сама функция для исполнения
			\item \textcolor{structure.fg}{pcName} - имя задачи (просто строка)
			\item \textcolor{structure.fg}{usStackDepth} - \textcolor{orange}{глубина стека задачи} (об этом позже)
			\item \textcolor{structure.fg}{pvParameters} - объект для передачи параметров если ничего не передаем - NULL
			\item \textcolor{structure.fg}{uxPriority} - приоритет задачи для планировщика - больше $\Leftrightarrow$ важнее
			\item \textcolor{structure.fg}{pxCreatedTask} - структура для управления задачей извне, можно передать NULL
			\item Отметим, что в FREERTOS \textcolor{orange}{есть одна всегда запущенная задача с низшим приоритетом - Idle Task} (простой) (она даже есть выше на диаграмме what actually happens).
			\item \textcolor{red}{Нельзя использовать RETURN!}
			\item \textcolor{red}{А как переключать задачи?}
	 \end{itemize}};
 \end{tikzpicture}
 \begin{frame}
 \large \textbf{\textcolor{structure.fg}{Запуск задачи}}\\
 \textcolor{red}{В файле main.cpp}:
 \includegraphics[width=10cm, height=3cm]{task6.png}\\
 \begin{tikzpicture}[remember picture,overlay]
 	
 	\node[anchor=north west] at ([xshift=0cm,yshift=-5.25cm]current page.north west) {
 		\begin{tabular}{c}
 			\includegraphics[width=5cm, height=3cm]{task5.png}\\
 			{main.cpp}
 		\end{tabular}
 	};
 \end{tikzpicture}
 \end{frame}
 \begin{frame}
 	
 
 \large \textbf{\textcolor{structure.fg}{Пример. Светодиод + аналоговый сигнал.}} \textcolor{red}{Код на github}\\
 \includegraphics[width=5cm, height=5cm]{example_1.png}
 \begin{tikzpicture}[remember picture,overlay]
 	\node[draw,rounded corners,fill=green!3,text width=6cm,align=left]
 	at (1,3) 
 	{\begin{itemize}
 			\item \textcolor{structure.fg}{Импорт библиотек}
 			\item \textcolor{structure.fg}{Далее прототипы функций}
 			\item \textcolor{structure.fg}{ADCValue} - где будем хранить данные с ADC. Не заубдьте про модификатор \textcolor{orange}{volatile}
 	\end{itemize}};
 \end{tikzpicture}
\end{frame}
\begin{frame}
 
 \large \textbf{\textcolor{structure.fg}{Пример. Запуск и инициализация.}} \textcolor{red}{Код на github}\\
 \includegraphics[width=12cm, height=6cm]{example2.png}
 \begin{tikzpicture}[remember picture,overlay]
 	\node[draw,rounded corners,fill=green!3,text width=8cm,align=left]
 	at (11,-0.3) 
 	{\begin{itemize}
 			\item \textcolor{structure.fg}{Инициализация периферии}
 			\item \textcolor{structure.fg}{Создание двух задач приоритета 1 (минимальный - 0 (Idle task))}
 			\item \textcolor{structure.fg}{Запуск планировщика задач} 
 	\end{itemize}};
 \end{tikzpicture}
\end{frame}
\begin{frame}
 
 \large \textbf{\textcolor{structure.fg}{Пример. Задача светодиода.}} \textcolor{red}{Код на github}\\
 \includegraphics[width=10cm, height=6cm]{example3.png}
 \begin{tikzpicture}[remember picture,overlay]
 	\node[draw,rounded corners,fill=green!3,text width=6cm,align=left]
 	at (1.5,3) 
 	{\begin{itemize}
 			\item \textcolor{structure.fg}{Изменение сигнала светодиода }1->0 или 0->1
 			\item \textcolor{structure.fg}{Задержка для мигания }500 тиков, т.е. 500мс
 	\end{itemize}};
 \end{tikzpicture}
 
\end{frame}
\begin{frame}
 \large \textbf{\textcolor{structure.fg}{Пример. Задача чтения аналогового сигнала}} \textcolor{red}{Код на github}\\
 \includegraphics[width=10cm, height=6.5cm]{example4.png}
 \begin{tikzpicture}[remember picture,overlay]
 	\node[draw,rounded corners,fill=green!3,text width=8cm,align=left]
 	at (11.3,0.75)  
 	{\begin{itemize}
 			\item \textcolor{structure.fg}{Ожидание готовности ADC}
 			\item \textcolor{structure.fg}{Чтение ADC}
 			\item \textcolor{structure.fg}{Опциональная задержка} - можете её не использовать
 			\item \textcolor{orange}{Не забудьте запустить ADC в int main()}\\ (ADC\_SoftwareStartConv(m\_ADCx))
 	\end{itemize}};
 \end{tikzpicture}
\end{frame}
\begin{frame}
 
 \large \textbf{\textcolor{structure.fg}{Синхронизации}}\\
 $\bullet$ Сейчас мы запускали задачи, которые работали независимо друг от друга. В реальных проектах - это редкость.\\
 $\bullet$ Можно выделить 3 основных вида взаимодействия задач:
 \begin{itemize}
 	\item Обмен данными
 	\item Ожидание других задач
 	\item Ограничения доступа к ресурсам
 \end{itemize}
\includegraphics[width=7.5cm, height=4.5cm]{mem2.jpg}
\end{frame}
\begin{frame}
 
  \large \textbf{\textcolor{structure.fg}{Обмен данными. Уведомления}}\\
  \includegraphics[width=7.5cm, height=4cm]{exchange1.png}
  \begin{itemize}
  	\item xHandler - как раз структура для контроля задачи извне (последний аргумент при запуске задачи).
  	\item xTaskNotifyTake - ожидание уведомление. Если уведомлений  нет - задача засыпает. В случае pdTrue - при выходе из TaskNotifyTake счетчик уведомлений обнуляется и задача проходит один раз. В случае pdFALSE - счетчик уменьшается на 1.
  	\item xTaskNotifyGive - уведомить задачу по её TaskHandler, т.е. увеличить счётчик уведомлений на 1.
  \end{itemize}
  \end{frame}
  \begin{frame}
  	\begin{itemize}
  	\item Для счетчика уведомления используют поле в структуре TCB (Task Control Block) - специальный блок, который выделяется после создания задачи.
  	\item В нем лежат uint32\_t ulNotifiedValue - ячейка с любым числом и uint8\_t ucNotifyState - байт состояния (ждем уведомления; не ждем; получили, но еще не прочитали).
  	\item Их и использует механизм уведомлений. Теоретически - в ячейку ulNotifiedValue можно положить любые данные, что позволяет существенно сэкономить память, не создавая очереди и т.п.
  \end{itemize}
\end{frame}
\begin{frame}[t]
	\large \textbf{\textcolor{structure.fg}{Обмен данными. StreamBuffer.} \textcolor{orange}{Отправка}}. Пример в \textcolor{red}{github}\\
	\begin{itemize}
		\item \textcolor{orange}{Stream Buffer} - это блокирующий способ обмена информацией между задачами. Все данные хранятся \textcolor{red}{побайтово}.
		\item Если в нем закончится место - он остановит задачу, которая попытается в него что-то положить, пока не пройдет TicksToWait или место не освободиться
		\item Очень удобен для работы с IO, например в случае UART у нас в любом случае приходят и отправляются байты, поэтому такой примитивной и \textcolor{orange}{легковесной структуры}нам хватит
	\end{itemize}
	\end{frame}
\begin{frame}[t]
  \large \textbf{\textcolor{structure.fg}{Обмен данными. StreamBuffer.} \textcolor{orange}{Отправка}}. Пример в \textcolor{red}{github}\\
  \includegraphics[width=6cm, height=4cm]{buffer1.png}\\
  \begin{tikzpicture}[remember picture,overlay]
  	\node[draw,rounded corners,fill=green!3,text width=8.5cm,align=left]
  	at (10.75,1.35) 
  	{\begin{itemize}
  			\item \textcolor{structure.fg}{xStreamBuffer} - хэндлер буфера, структура аналогичная хэндлеру для задачи, т.е. а.к.а. отссылка на конкретный буффер в памяти 
  			\item \textcolor{structure.fg}{*pvTxDara} - указатель на начало фрагмента памяти с данными
  			\item \textcolor{structure.fg}{xDataLengthBytes} - длина фрагмента памяти в байтах
  			\item \textcolor{structure.fg}{xTicksToWait} - тайм-аут если в буфере закончилось место (Защита от вечной блокировки).
  	\end{itemize}};
  \end{tikzpicture}
  
\end{frame}
\begin{frame}[t]
  
  \large \textbf{\textcolor{structure.fg}{Обмен данными. StreamBuffer.}\textcolor{orange}{Приём}}. Пример в \textcolor{red}{github}\\
  \includegraphics[width=6cm, height=4cm]{buffer2.png}\\
  \begin{tikzpicture}[remember picture,overlay]
  	\node[draw,rounded corners,fill=green!3,text width=9cm,align=left]
  	at (10.75,1.5) 
  	{\begin{itemize}
  			\item \textcolor{structure.fg}{xStreamBuffer} - хэндлер буфера, структура аналогичная хэндлеру для задачи, т.е. а.к.а. отссылка на конкретный буффер в памяти 
  			\item \textcolor{structure.fg}{*pvRxDara} - указатель на начало фрагмента памяти для чтения
  			\item \textcolor{structure.fg}{xDataLengthBytes} - таймаут блокировки, если буфер пустой (нечего читать).
  			\item \textcolor{structure.fg}{xTicksToWait} - тайм-аут если в буфере закончилось место (Защита от вечной блокировки).
  			\item textcolor{red}{Trigger Level} - необязательный параметр, который устанавливает минимальное количество байт в буфере чтобы проснуться (например для общения по UART это нужно, чтобы вы могли ожидать прием всего пакета, а не нескольких байт из него) (Trigger level есть в примере на github).
  	\end{itemize}};
  \end{tikzpicture}
  
\end{frame}

\begin{frame}[t]
\large \textbf{\textcolor{structure.fg}{Обмен данными. Queue.}\textcolor{orange}{Приём}}. Пример в \textcolor{red}{github}\\
Следующий метод это очередь. Работает схожим образом, но по умолчанию хранит структуры в одной ячейке\\
  \includegraphics[width=6cm, height=3.5cm]{queue1.png} \includegraphics[width=6cm, height=3.5cm]{queue2.png}\\
  \begin{tikzpicture}[remember picture,overlay]
  	\node[draw,rounded corners,fill=green!3,text width=14cm,align=left]
  	at (6.75,-0.75) 
  	{\begin{itemize}
  			\item \textcolor{structure.fg}{xQueue} - хэндлер очереди, структура аналогичная хэндлеру для задачи, т.е. а.к.а. отссылка на конкретную очередь в памяти 
  			\item \textcolor{structure.fg}{*pvItemQueue} - указатель на объект куда/откуда писать 
  			\item \textcolor{structure.fg}{xTicksToWait} - тайм-аут если в буфере закончилось место (Защита от вечной блокировки).
  	\end{itemize}};
  \end{tikzpicture}
  \end{frame}
  \begin{frame}
  	\large \textbf{\textcolor{structure.fg}{Ожидание задач. Семафоры}}. Пример в \textcolor{red}{github}
  	\begin{itemize}
  		\item \textcolor{orange}{Семафор} - это счётчик разрешений. Можно делать бинарными (1 разрешение), можно множественными.
  		\item Задача берет разрешение, если оно есть, то счётчик уменьшается на 1, иначе зависает, в ожидании, что другая задача отдаст разрешение (в этом случае счетчик увеличиться на 1).
  		\item При неудаче (нет доступных разрешений и прошел тайм-аут возвращает какое-то значение, можете его погуглить)

  	\end{itemize}
  \end{frame}
  \begin{frame}
  		\large \textbf{\textcolor{structure.fg}{Ожидание задач. Виды семафоров}}. Пример в \textcolor{red}{github}
  		\begin{itemize}
  			\item Отличаются только тем, что вы проинициализируете для первого аргумента \textcolor{red}{(на гитхабе пример бинарного семафора)}
  			\item Бинарный семафор - одно разрешение
  			\item Счётный семафор - внутри счётчик разрешений
  			\item Mutex - Mutual Exclusion. Нужен для разграничения доступа к ресурсам. Это бинарный семафор с важным условием - возвращать его может только тот, кто его взял. \textcolor{orange}{Это концепция владения ресурсом}. Нельзя использовать в прерываниях $\Rightarrow$ используйте бинарный семафор.
  		\end{itemize}
  	\includegraphics[width=6cm, height=3.5cm]{sem1.png} \includegraphics[width=6cm, height=3.5cm]{sem2.png}\\
  	\end{frame}
  	\begin{frame}
  		\large \textbf{\textcolor{structure.fg}{Группы событий}}. Пример в \textcolor{red}{github}.
  		\begin{itemize}
  			\item \textcolor{orange}{EventGroups} - механизм для ожидания множественных флагов
  			\item Какие-то задачи выставляют флаги какие-то флаги, сигнализируя события (например блок работы с камерой сигнализирует прием новых данных, а блок управления кикером сигнализирует, что удар перезарядился).
  			\item Другие задачи могут засыпать в ожидании определенных комбинаций битов или хотя бы одного
  		\end{itemize}
  		
  	\end{frame}
  	\begin{frame}[t]
  			\large \textbf{\textcolor{structure.fg}{Группы событий}}. Пример в \textcolor{red}{github}.\\
  		\includegraphics[width=6cm, height=3.5cm]{groups1.png} \includegraphics[width=6cm, height=3.5cm]{groups2.png}\\
  		\begin{tikzpicture}[remember picture,overlay]
  			\node[draw,rounded corners,fill=green!3,text width=14cm,align=left]
  			at (6.75,-1.75) 
  			{\begin{itemize}
  					\item \textcolor{structure.fg}{xEventGroup} - хэндлер группы.
  					\item \textcolor{structure.fg}{uxBitsToSet} - какие биты выставить.
  					\item \textcolor{structure.fg}{uxBitsWaitFor} - Какие биты ждать.
  					\item \textcolor{structure.fg}{xClearOnExit} - При срабатывании группы: pdTrue (аналог true для FreeRtos) $\implies$ сбросить биты, иначе (pdFalse) не сбрасывать.
  					\item \textcolor{structure.fg}{xWaitForAllBits} - pdTrue - ждать все биты (логическое И), иначе ждать хотя бы один (логическое или).
  					\item \textcolor{structure.fg}{xTicksToWait} - тайм-аут.
  			\end{itemize}};
  		\end{tikzpicture}
  	\end{frame}
  	
  	\begin{frame}
  		\large \textbf{\textcolor{structure.fg}{Взаимодействие задач в прерываниях}}.
  		\begin{itemize}
  			\item Прерывание - это не обычная задача. 
  			\item Она не учитывается планировщиком, у нее нет стека, хэндера и т.п.
  			\item Блокировка в прерывании вызовет зависание программы (т.к. системных тиков нет)
  			\item Поэтому все вышеописанные механизмы не будут работать в прерываниях
  			\item Для всех структур (кроме mutex -> вместо него бинарный семафор) есть fromISR альтернативы, которые дают не прямую директиву планировщику, а указание, что делать сразу после выхода из прерывания.
  			\item Если всё же внутри прерывания нужен контроль ресурсов, то используются критические секции (мы их уже использовали для UART).
  			\newpage
  			\item Везде нужно просто в конце добавлять fromISR, в случае очереди: xQueueSend() $\implies$ xQueueSend\textcolor{red}{FromISR}()
  			\item Также для каждой функции появляется доп аргумент, который указывает планировщику, нужно ли переключать контекст при выходе из прерывания, или идем по стандартному пути (откуда вошли в прерывание туда и вернулись, или переключились на доп действие).
  			\item В конце прерывания нужно вызвать portYIELD\_FROM\_ISR(xHigherPriorityTaskWoken); Оно как раз и вызовет внеплановое переключение контекста.
  			\item \textcolor{red}{Пример прерывания есть в github!}
  		\end{itemize}
  	\end{frame}
\end{document}
 

